%------------------------------------------------------------------------------%
\begin{question}
  Considerando a equação
  \[
    \frac{(x-1)^2}{4} + \frac{(y+2)^2}{9} = 1
  \]
  \vspace{-\baselineskip}
  \begin{enumerate}
    \item identifique a curva
    \item determine seus principais valores e pontos
    \item esboce a curva
  \end{enumerate}
\end{question}

%------------------------------------------------------------------------------%
\begin{resolution}{Identificação}
\begin{columns}[t]
\column{0.4\textwidth}
  \onslide<+->{A equação
  \[
    \frac{(x-1)^2}{4} + \frac{(y+2)^2}{9} = 1
  \]}%
  \onslide<+->{está na forma de uma elipse transladada
  \[
   \frac{(x-h)^2}{a^2} + \frac{(y-k)^2}{b^2} = 1
  \]}
  \onslide<+->{Identificamos o centro}
  \[
    \onslide<+->{C = (h,  k)}
    \onslide<+->{  = (1, -2)}
  \]

\column{0.4\textwidth}
  \onslide<+->{Os semieixos}
  \begin{align*}
    \onslide<+->{a^2 & = 4}
    \qquad
    \onslide<+->{a = 2}
    \\
    \onslide<+->{b^2 & = 9}
    \qquad
    \onslide<+->{b = 3}
  \end{align*}
  \onslide<+->{Vértices}
  \begin{align*}
    \onslide<+->{V_1
    & = C + (a, 0)}
    \onslide<+->{= (1 + 2, -2)}
    \onslide<+->{= (3, -2)}
    \\
    \onslide<+->{V_2
    & = C - (a, 0)}
    \onslide<+->{= (1-2, -2)}
    \onslide<+->{= (-1, 1)}
    \\
    \onslide<+->{V_3
    & = C + (0, b)}
    \onslide<+->{= (1, -2 - 3)}
    \onslide<+->{= (1, -5)}
    \\
    \onslide<+->{V_4
    & = C - (0, b)}
    \onslide<+->{= (1, -2 + 3)}
    \onslide<+->{= (1, 1)}
  \end{align*}
\end{columns}
\end{resolution}

%------------------------------------------------------------------------------%
\begin{resolution}{Esboço}
  \centering
  \includegraphics{I-07-Identificacao-ex-7}
\end{resolution}

%------------------------------------------------------------------------------%
